\section{Architektur \& Befehle}

\subsection{ISA Thumb-2 Befehlssatz}
\begin{itemize}
    \item Verarbeitet 32 und 16 Bit Befehle, Code-Density, Performance, mehr Befehle möglich.
    \item Flexibler 2. Operand $<$op2$>$
    \item Erweiterte Addressing-Modes
    \item Erweiterte Grundoperationen (z.B. OR NOT, für jeden 16Bit-Befehl neu auch einen 32Bit-Befehl)
    \item Bit-Field Processing
    \item Compare \& Branch Zero, Conditional Execution (IT-Instruction)
    \item Saturation (32Bit)
    \item Multiply \& Accumulate (32Bit MAC)
\end{itemize}

\subsubsection{Single Instruction Multiple Data (SIMD)}

\begin{itemize}
    \item ganze Arrays von Daten werden übergeben.
    \item Saturation (8/16Bit)
    \item Multiply \& Add/Sub oder Accumulate (8/16Bit)
    \item Add \& Sub (8/16Bit)
    \item Add/Sub \& Saturation
\end{itemize}

\subsubsection{Load/Store-Architektur (RISC)}
Sign Extender, Barrel-Shifter (Erlaubt Shift-Operationen auf R).

\subsubsection{Multiply \& Accumulate (MAC)}
Erlaubt Rechenoperationen mit 3 Werten.

\subsection{Adressing Modes}
\begin{ctabular}{@{}ll@{}}
\toprule
\textbf{Mode} & \textbf{Beispiel} \\
\midrule
Direkt & \texttt{R1} or \texttt{\#100} \\
Liste & \texttt{(R1-R3, R5, LR)} \\
Indirekt & \texttt{[R0]} \\
Indirekt Offset & \texttt{[R0, \#4]} \\
Indirekt Index & \texttt{[R0, R1]} \\
Indirekt Pre-Index & \texttt{[R0, \#2]!} \\
Indirekt Post-Index & \texttt{[R0], \#-8} \\
Indirekt Shifted Index & \texttt{[R0, R1, LSL \#1]} \\
PC-relative & \texttt{label1} or \texttt{[PC, \#4]} \\
SP-relativ & \texttt{[SP, \#4]} \\
\bottomrule
\end{ctabular}

\vspace{0.5em}
\noindent\framebox{\parbox{\linewidth}{\centering \vspace{10pt}\textit{[Placeholder: Architektur Diagramm (MAC, RO-R15, Data Bus, Address Bus)]}\vspace{10pt}}}
\vspace{0.5em}
